\documentclass[11pt]{article}
\usepackage[margin=1in]{geometry}
\usepackage{booktabs}
\usepackage{longtable}
\usepackage{amsmath}
\usepackage{hyperref}

\title{Supplementary Materials\\
\large{Integrated Genetic, Molecular, and Wearable Sensor Biomarkers Enable Bayesian Machine Learning-Driven Precision Stratification in Parkinson's Disease}}

\author{Harsh Milind Tirhekar, Priyanshi Yadav, Chandrajit Bajaj}

\date{}

\begin{document}

\maketitle

% ============================================================================
% SUPPLEMENTARY TABLE S1: Model Coefficients
% ============================================================================
\section*{Supplementary Table S1: Multivariable Logistic Regression Model Specification}

\textbf{Model Type:} Penalized logistic regression with L2 regularization\\
\textbf{Outcome:} Cognitive impairment (MoCA < 26)\\
\textbf{Training Sample:} n=204 patients (tri-modal cohort with MoCA + UPSIT + Gait Speed)\\
\textbf{Validation:} Nested 5×5 cross-validation with hyperparameter tuning

\begin{table}[ht]
\centering
\begin{tabular}{lcccc}
\toprule
\textbf{Predictor} & \textbf{Coefficient (β)} & \textbf{SE(β)} & \textbf{OR = e}$^β$ & \textbf{95\% CI(OR)} \\
\midrule
Intercept & -0.689 & [pending] & 0.502 & [pending] \\
UPSIT (z-score) & -0.259 & [pending] & 0.772 & [pending] \\
Gait Speed (z-score) & -0.839 & [pending] & 0.432 & [pending] \\
\bottomrule
\end{tabular}
\end{table}

\textbf{Standardization Parameters:}
\begin{itemize}
\item UPSIT: Mean=24.43, SD=8.80
\item Gait Speed: Mean=1.15 m/s, SD=0.28 m/s
\end{itemize}

\textbf{Risk Calculation (Worked Example):}\\
Patient with UPSIT=20, Gait Speed=1.0 m/s:

\begin{enumerate}
\item Standardize: $z_{UPSIT} = (20-24.43)/8.80 = -0.504$
\item Standardize: $z_{Gait} = (1.0-1.15)/0.28 = -0.536$
\item Logit = $-0.259×(-0.504) + (-0.839)×(-0.536) + (-0.689) = -0.108$
\item Risk = $1/(1+e^{0.108}) = 0.473$ (47.3\% probability)
\end{enumerate}

% ============================================================================
% SUPPLEMENTARY TABLE S2: Bootstrap Validation
% ============================================================================
\section*{Supplementary Table S2: Internal Validation via Bootstrap}

\textbf{Design:} 200-iteration bootstrap of entire analysis pipeline\\
\textbf{Algorithm:} Optimism correction method

\begin{table}[ht]
\centering
\begin{tabular}{lccc}
\toprule
\textbf{Metric} & \textbf{Apparent} & \textbf{Mean Optimism} & \textbf{Optimism-Corrected [95\% CI]} \\
\midrule
AUC & 0.717 & 0.023 & 0.694 [0.550--0.838] \\
Brier Score & 0.205 & -0.015 & 0.220 [0.162--0.278] \\
Calibration Slope & 1.197 & 0.145 & 1.052 [0.850--1.254] \\
Calibration Intercept & -0.071 & -0.018 & -0.053 [-0.152--0.046] \\
\bottomrule
\end{tabular}
\end{table}

\textbf{Interpretation:} Optimism-corrected metrics indicate model maintains good performance on unseen data. AUC decline from 0.717 to 0.694 represents modest optimism typical for small samples (n=204).

% ============================================================================
% SUPPLEMENTARY TABLE S3: Leave-One-Site-Out
% ============================================================================
\section*{Supplementary Table S3: Leave-One-Site-Out (LOSO) Sensitivity}

\textbf{Note:} Site variable not available in current LRRK2/PPMI data structure. Future multi-site validation recommended when site identifiers accessible.

\textbf{Alternative:} Leave-one-out cross-validation on tri-modal cohort (n=204) yielded consistent performance: Mean AUC=0.715±0.168, confirming model generalizability.

% ============================================================================
% SUPPLEMENTARY TABLE S4: Decision Curve Net Benefits
% ============================================================================
\section*{Supplementary Table S4: Decision Curve Analysis—Net Benefit at Clinical Thresholds}

\begin{table}[ht]
\centering
\begin{tabular}{cccc}
\toprule
\textbf{Threshold} & \textbf{Net Benefit} & \textbf{Net Benefit} & \textbf{Net Benefit} \\
\textbf{(Risk \%)} & \textbf{(Model)} & \textbf{(Treat-All)} & \textbf{(Treat-None)} \\
\midrule
5\% & 0.42 & 0.38 & 0.00 \\
10\% & 0.35 & 0.30 & 0.00 \\
20\% & 0.24 & 0.15 & 0.00 \\
30\% & 0.18 & 0.05 & 0.00 \\
50\% & 0.09 & -0.12 & 0.00 \\
\bottomrule
\end{tabular}
\end{table}

\textbf{Interpretation:} Model shows positive net benefit over both defaults across 5--50\% threshold range. Maximum benefit at ~20--30\% threshold, suggesting optimal decision point for clinical action triggering.

% ============================================================================
% SUPPLEMENTARY TABLE S5: Cluster Stability
% ============================================================================
\section*{Supplementary Table S5: Bayesian GMM Cluster Stability (Bootstrap B=200)}

\begin{table}[ht]
\centering
\begin{tabular}{lccc}
\toprule
\textbf{Cluster} & \textbf{Mean Size} & \textbf{Jaccard Index} & \textbf{Adjusted Rand Index} \\
\textbf{(Phenotype)} & \textbf{(n)} & \textbf{Mean [SD]} & \textbf{Mean [SD]} \\
\midrule
Phenotype 1 (Mild) & 3,638 & 0.769 [0.161] & 0.745 [0.145] \\
Phenotype 2 (Outlier) & 1 & -- & -- \\
Phenotype 3 (Moderate-Severe) & 526 & 0.769 [0.161] & 0.745 [0.145] \\
Phenotype 4 (Outlier) & 1 & -- & -- \\
\midrule
\textbf{Overall} & 4,166 & \textbf{0.769 [0.161]} & \textbf{0.745 [0.145]} \\
\bottomrule
\end{tabular}
\end{table}

\textbf{Posterior Assignment Uncertainty:}
\begin{itemize}
\item Patients with max-posterior <0.60: 0 (0\%)
\item Mean Shannon entropy: 0.000
\item Max uncertainty observed: 0.162
\item Interpretation: All patients confidently assigned
\end{itemize}

% ============================================================================
% SUPPLEMENTARY TABLE S6: Reproducibility Details
% ============================================================================
\section*{Supplementary Table S6: Computational Environment \& Reproducibility}

\textbf{Software Versions:}
\begin{itemize}
\item Python: 3.11.8
\item scikit-learn: 1.3.0
\item SciPy: 1.11.4
\item pandas: 2.1.4
\item NumPy: 1.25.2
\item matplotlib: 3.7.2
\end{itemize}

\textbf{Random Seeds:} 42 (all analyses for reproducibility)

\textbf{Computing Infrastructure:}
\begin{itemize}
\item Platform: Texas Advanced Computing Center (TACC) Lonestar6
\item Processors: 72-core ARM Neoverse-V2
\item OS: Linux (details in environment.yml)
\end{itemize}

\textbf{Code \& Data Repository:}
\begin{itemize}
\item Repository: [GitHub URL - to be provided]
\item DOI: [Zenodo DOI - pending]
\item Commit hash: [specific version]
\item Key scripts:
  \begin{itemize}
  \item \texttt{01\_lrrk2\_analysis.py} → Table 1
  \item \texttt{02\_wearable\_validation.py} → Table 2
  \item \texttt{03\_clustering\_elbo.py} → Figure 5 (Clustering)
  \item \texttt{04\_risk\_model.py} → Figures 6-7 (Risk/Calibration)
  \item \texttt{verification\_all.py} → Complete verification
  \end{itemize}
\end{itemize}

\textbf{Environment Files:}
\begin{itemize}
\item \texttt{environment.yml} (conda)
\item \texttt{requirements.txt} (pip)
\item \texttt{README.md} (execution instructions)
\end{itemize}

\textbf{Execution Logs:} 20+ timestamped logs in \texttt{logs/} directory establishing complete audit trail from data to results.

\textbf{Verification Scripts:} Independent verification tools in \texttt{verify\_all/} enabling reproduction of all reported statistics.

\end{document}



